\section{Çalışma Alanı}

Bu çalışmada Hakkari ili ve çevresi incelenmiştir. Bölge, Türkiye'nin güneydoğusunda yer almakta olup yüksek rakımlı dağlık bir yapıya sahiptir.

\subsection{Meteoroloji İstasyonları}

Çalışmada kullanılan meteoroloji istasyonları Tablo \ref{tab:istasyonlar}'da listelenmiştir.

\begin{table}[H]
\centering
\caption{Çalışmada Kullanılan Meteoroloji İstasyonları}
\label{tab:istasyonlar}
\begin{tabular}{llrrr}
\toprule
İstasyon No & İstasyon Adı & Enlem & Boylam & Rakım (m) \\
\midrule
17285 & Hakkari Merkez & 37.5745 & 43.7388 & 1727 \\
17920 & Yüksekova & 37.5785 & 44.2862 & 1877 \\
18234 & Şemdinli & 37.2950 & 44.5819 & 1284 \\
18771 & Durankaya & 37.5600 & 43.6094 & 1825 \\
17815 & Yüksekova Havalimanı & 37.5481 & 44.2399 & 1854 \\
19909 & Kayak Merkezi & 37.5714 & 43.6818 & 2516 \\
\bottomrule
\end{tabular}
\end{table}

\section{Veri Kaynakları}

\subsection{Yıldırım Verileri}

Yıldırım verileri, Meteoroloji Genel Müdürlüğü Yıldırım Tespit Sistemi'nden (YILTES) elde edilmiştir. Veri seti aşağıdaki bilgileri içermektedir:

\begin{itemize}
    \item Olay zamanı (yıl-ay-gün saat:dakika:saniye.milisaniye)
    \item Konum (enlem, boylam - WGS84 koordinat sistemi)
    \item Yükseklik (km)
    \item Akım miktarı (kA)
    \item Olay tipi (Yıldırım/Şimşek)
    \item Referans noktasına mesafe (km)
\end{itemize}

Toplamda yaklaşık 258,000 yıldırım/şimşek olayı kaydedilmiştir (2015-2025).

\subsection{Meteorolojik Veriler}

Günlük meteorolojik veriler Meteoroloji Genel Müdürlüğü'nden temin edilmiştir:

\begin{itemize}
    \item Maksimum, minimum ve ortalama sıcaklık (°C)
    \item Ortalama aktüel basınç (hPa)
    \item Ortalama nispi nem (\%)
    \item Ortalama bulutluluk (8 Okta)
    \item Maksimum rüzgar hızı ve yönü (m/s)
    \item Toplam yağış (mm)
    \item Yatay görüş mesafesi (km)
\end{itemize}

\section{Veri Ön İşleme}

\subsection{Veri Temizleme}

\begin{enumerate}
    \item Tarih formatlarının standardizasyonu
    \item Eksik değerlerin incelenmesi
    \item Aykırı değerlerin tespiti
\end{enumerate}

\subsection{Konum Eşleştirmesi}

Her yıldırım olayı, koordinatlarına en yakın meteoroloji istasyonuna eşleştirilmiştir. Bu işlem Öklid mesafesi kullanılarak gerçekleştirilmiştir.

\subsection{Günlük Agregasyon}

Yıldırım verileri istasyon ve tarih bazında günlük olarak özetlenmiştir:
\begin{itemize}
    \item Günlük yıldırım sayısı
    \item Ortalama akım değeri
    \item Maksimum akım değeri
\end{itemize}

\section{Veri Seti İstatistikleri}

Şekil \ref{fig:yillik} ve Şekil \ref{fig:aylik} yıldırım olaylarının zamansal dağılımını göstermektedir.

\begin{figure}[H]
\centering
\includegraphics[width=0.9\textwidth]{yillik_yildirim.png}
\caption{Yıllara Göre Yıldırım/Şimşek Olayları}
\label{fig:yillik}
\end{figure}

\begin{figure}[H]
\centering
\includegraphics[width=0.9\textwidth]{aylik_yildirim.png}
\caption{Aylara Göre Yıldırım/Şimşek Dağılımı}
\label{fig:aylik}
\end{figure}

\section{Hedef Değişken}

Hedef değişken olarak ikili sınıflandırma kullanılmıştır:
\begin{equation}
    y = \begin{cases} 
        1 & \text{günde en az 1 yıldırım olayı var} \\
        0 & \text{yıldırım olayı yok}
    \end{cases}
\end{equation}
